\chapter{The ladder}
\label{ch:ladder}

The ladder is the list of relations that the compiler knows before it meets a machine. ladder.h states 21
anchors and 109 cases. Each case is an anchor, one or two words of 32 places, and the word that the relation
gives them. Lean reads these cases as data that the tool \texttt{harness.py lng prove} writes into
\texttt{Lng/ladder\_cases.lean}.

\section{Words and relations}

\begin{definition}
A word of $n$ places is an element of $W_n = \mathbb{Z}/2^n\mathbb{Z}$, the type \texttt{BitVec n} of Lean. The
truth of a predicate $p$ on $W_n$ is the word $\mathbf{1}_n$ of ones where $p$ holds, and $0$ where it fails.
\end{definition}

\begin{definition}
For a sign reading $s \in \{\mathrm{u}, \mathrm{s}\}$, the relation $R^{s}_{a}$ of an anchor $a$ is the partial
map on $W_n \times W_n$ that \texttt{relation} defines in \texttt{Lng/relation.lean}. The orderings compare as
unsigned words under $\mathrm{u}$ and as two's complement words under $\mathrm{s}$. Under $\mathrm{s}$ a
quotient truncates toward 0, and a remainder takes the sign of the dividend, as C11 states
\cite{src:c11}. DOWN under $\mathrm{s}$ fills the emptied places with the sign. A quotient or a remainder by 0 has
no value. UP and DOWN read their count modulo $n$.
\end{definition}

The signed reading follows \texttt{ladder\_signed\_answer} of ladder.h. Its one edge case is the lowest word over a
word of ones, where the true quotient $2^{n-1}$ has no place in $W_n$.

\begin{lemma}[\texttt{lowest\_over\_ones}]
At 32 places under $\mathrm{s}$, the quotient of $\mathtt{0x80000000}$ by $\mathtt{0xffffffff}$ is
$\mathtt{0x80000000}$, and the remainder is 0.
\end{lemma}

\section{The cases}

The ladder law of \texttt{Lng/ladder\_law.lean} gives the answer of each anchor on words of 32 places. It is the
unsigned relation, read on natural numbers.

\begin{theorem}[\texttt{cases\_answered}]
For each case $(a, w, y)$ of ladder.h, the ladder law gives the anchor $a$ the answer $y$ on the words $w$.
\end{theorem}

\begin{proof}
The statement is a conjunction over 109 closed terms. The kernel decides it.
\end{proof}

\begin{theorem}[\texttt{agrees}]
For every anchor text $a$ and all $x, y \in W_{32}$, the value of $R^{\mathrm{u}}_a(x, y)$ is the answer of the
ladder law for $a$ on $x$ and $y$.
\end{theorem}

\begin{proof}
Both sides branch on the same anchor texts. Each branch reduces by the definitions of the words and of their
natural values. The two branches of a quotient and of a remainder split again on a divisor of 0.
\end{proof}

The two theorems close the gap between the general relation and the cases. A relation at any width reduces, at
32 places and unsigned, to the law that every case of ladder.h obeys.
